%%%%%%%%%%%%%%%%%%%%%%%%%%%%%%%%%%%%%%%%%%%%%%%%%%%%%%%%%%%%
% 5.9 FUNCTIONAL ANALYSIS
% The Composition of Advanced Mathematics
%%%%%%%%%%%%%%%%%%%%%%%%%%%%%%%%%%%%%%%%%%%%%%%%%%%%%%%%%%%%

\section{Functional Analysis}
\label{sec:functional-analysis}

\prerequisite{
Linear algebra, metric spaces, topology, real analysis, measure theory,
Lebesgue integration, sequences and series of functions, and basic complex
analysis.
}

\learningobjectives{
By the end of this section, the reader should be able to:
\begin{itemize}
    \item explain the purpose of functional analysis;
    \item work with normed vector spaces and Banach spaces;
    \item compare equivalent norms;
    \item identify important examples such as \(\ell^p\), \(L^p\), and \(C(K)\);
    \item determine whether linear operators are bounded and continuous;
    \item compute and interpret operator norms;
    \item understand continuous linear functionals and dual spaces;
    \item apply the Hahn--Banach Theorem;
    \item understand the Uniform Boundedness Principle;
    \item apply the Open Mapping and Closed Graph Theorems;
    \item work with inner-product spaces and Hilbert spaces;
    \item use orthogonality and orthogonal projections;
    \item understand orthonormal sets and bases;
    \item apply Bessel's inequality and Parseval's identity;
    \item use the Riesz Representation Theorem for Hilbert spaces;
    \item distinguish strong and weak convergence;
    \item understand compact operators;
    \item define the spectrum and resolvent of an operator;
    \item recognize the basic spectral theory of bounded and
          self-adjoint operators.
\end{itemize}
}

%%%%%%%%%%%%%%%%%%%%%%%%%%%%%%%%%%%%%%%%%%%%%%%%%%%%%%%%%%%%
\subsection{What Is Functional Analysis?}
\label{subsec:what-is-functional-analysis}
%%%%%%%%%%%%%%%%%%%%%%%%%%%%%%%%%%%%%%%%%%%%%%%%%%%%%%%%%%%%

Linear algebra studies vectors and linear transformations, usually in
finite-dimensional spaces such as

\[
\R^n
\]

or

\[
\mathbb{C}^n.
\]

Analysis studies limits, convergence, continuity, differentiation, and
integration.

Functional analysis combines these viewpoints by studying
infinite-dimensional vector spaces equipped with analytic structure.

Typical vectors are no longer finite coordinate lists. They may be

\[
\boxed{
\text{functions, sequences, signals, or solutions of differential equations}.
}
\]

Typical spaces include

\[
C([a,b]),
\qquad
L^p(X,\mu),
\qquad
\ell^p,
\]

and Sobolev spaces.

The central objects are often linear operators

\[
\boxed{
T:X\to Y
}
\]

between such spaces.

%%%%%%%%%%%%%%%%%%%%%%%%%%%%%%%%%%%%%%%%%%%%%%%%%%%%%%%%%%%%
\subsection{Normed Vector Spaces}
\label{subsec:normed-vector-spaces}
%%%%%%%%%%%%%%%%%%%%%%%%%%%%%%%%%%%%%%%%%%%%%%%%%%%%%%%%%%%%

\begin{definitionbox}{Norm}
Let \(X\) be a vector space over \(\R\) or \(\mathbb{C}\).

A norm is a function

\[
\|\cdot\|:X\to[0,\infty)
\]

satisfying, for all \(x,y\in X\) and scalars \(\alpha\),

\[
\boxed{
\|x\|\geq0,
}
\]

\[
\boxed{
\|x\|=0
\iff
x=0,
}
\]

\[
\boxed{
\|\alpha x\|
=
|\alpha|\|x\|,
}
\]

and

\[
\boxed{
\|x+y\|
\leq
\|x\|+\|y\|.
}
\]
\end{definitionbox}

A vector space equipped with a norm is called a normed vector space.

%%%%%%%%%%%%%%%%%%%%%%%%%%%%%%%%%%%%%%%%%%%%%%%%%%%%%%%%%%%%
\subsection{Norms Produce Metrics}
\label{subsec:norms-produce-metrics}
%%%%%%%%%%%%%%%%%%%%%%%%%%%%%%%%%%%%%%%%%%%%%%%%%%%%%%%%%%%%

Every norm determines a metric by

\[
\boxed{
d(x,y)=\|x-y\|.
}
\]

Therefore every normed space is automatically a metric space.

Consequently, concepts such as

\[
\boxed{
\text{open sets, closed sets, convergence, Cauchy sequences, and continuity}
}
\]

apply immediately.

A sequence \((x_n)\) converges to \(x\) in norm if

\[
\boxed{
\|x_n-x\|\to0.
}
\]

%%%%%%%%%%%%%%%%%%%%%%%%%%%%%%%%%%%%%%%%%%%%%%%%%%%%%%%%%%%%
\subsection{Examples of Norms on Finite-Dimensional Spaces}
\label{subsec:finite-dimensional-norms}
%%%%%%%%%%%%%%%%%%%%%%%%%%%%%%%%%%%%%%%%%%%%%%%%%%%%%%%%%%%%

For

\[
x=(x_1,\ldots,x_n)\in\R^n,
\]

important norms include

\[
\boxed{
\|x\|_1
=
\sum_{k=1}^{n}|x_k|,
}
\]

\[
\boxed{
\|x\|_2
=
\left(
\sum_{k=1}^{n}|x_k|^2
\right)^{1/2},
}
\]

and

\[
\boxed{
\|x\|_\infty
=
\max_{1\leq k\leq n}|x_k|.
}
\]

More generally, for \(1\leq p<\infty\),

\[
\boxed{
\|x\|_p
=
\left(
\sum_{k=1}^{n}|x_k|^p
\right)^{1/p}.
}
\]

%%%%%%%%%%%%%%%%%%%%%%%%%%%%%%%%%%%%%%%%%%%%%%%%%%%%%%%%%%%%
\subsection{Equivalent Norms}
\label{subsec:equivalent-norms}
%%%%%%%%%%%%%%%%%%%%%%%%%%%%%%%%%%%%%%%%%%%%%%%%%%%%%%%%%%%%

\begin{definitionbox}{Equivalent Norms}
Two norms \(\|\cdot\|_a\) and \(\|\cdot\|_b\) on \(X\) are equivalent if
there exist constants \(c,C>0\) such that

\[
\boxed{
c\|x\|_a
\leq
\|x\|_b
\leq
C\|x\|_a
}
\]

for every \(x\in X\).
\end{definitionbox}

Equivalent norms generate the same notion of convergence and the same
topology.

A fundamental finite-dimensional result is

\[
\boxed{
\text{all norms on a finite-dimensional vector space are equivalent}.
}
\]

This statement generally fails in infinite-dimensional spaces.

%%%%%%%%%%%%%%%%%%%%%%%%%%%%%%%%%%%%%%%%%%%%%%%%%%%%%%%%%%%%
\subsection{Sequence Spaces}
\label{subsec:sequence-spaces-functional}
%%%%%%%%%%%%%%%%%%%%%%%%%%%%%%%%%%%%%%%%%%%%%%%%%%%%%%%%%%%%

Functional analysis treats infinite sequences as vectors.

For

\[
1\leq p<\infty,
\]

define

\[
\boxed{
\ell^p
=
\left\{
x=(x_n):
\sum_{n=1}^{\infty}|x_n|^p<\infty
\right\}.
}
\]

The norm is

\[
\boxed{
\|x\|_p
=
\left(
\sum_{n=1}^{\infty}|x_n|^p
\right)^{1/p}.
}
\]

For \(p=\infty\),

\[
\boxed{
\ell^\infty
=
\left\{
x=(x_n):
\sup_n|x_n|<\infty
\right\},
}
\]

with

\[
\boxed{
\|x\|_\infty
=
\sup_n|x_n|.
}
\]

%%%%%%%%%%%%%%%%%%%%%%%%%%%%%%%%%%%%%%%%%%%%%%%%%%%%%%%%%%%%
\subsection{The Spaces \(c_0\) and \(c\)}
\label{subsec:c0-c-spaces}
%%%%%%%%%%%%%%%%%%%%%%%%%%%%%%%%%%%%%%%%%%%%%%%%%%%%%%%%%%%%

Two important subspaces of \(\ell^\infty\) are

\[
\boxed{
c
=
\{(x_n):x_n\text{ converges}\}
}
\]

and

\[
\boxed{
c_0
=
\{(x_n):x_n\to0\}.
}
\]

Both use the supremum norm

\[
\boxed{
\|x\|_\infty
=
\sup_n|x_n|.
}
\]

The space \(c_0\) is especially important in the study of dual spaces.

%%%%%%%%%%%%%%%%%%%%%%%%%%%%%%%%%%%%%%%%%%%%%%%%%%%%%%%%%%%%
\subsection{Spaces of Continuous Functions}
\label{subsec:continuous-function-spaces}
%%%%%%%%%%%%%%%%%%%%%%%%%%%%%%%%%%%%%%%%%%%%%%%%%%%%%%%%%%%%

Let \(K\) be a compact topological space.

The space

\[
\boxed{
C(K)
}
\]

consists of continuous scalar-valued functions on \(K\).

A natural norm is the supremum norm

\[
\boxed{
\|f\|_\infty
=
\sup_{x\in K}|f(x)|.
}
\]

For example,

\[
C([0,1])
\]

with this norm is one of the standard infinite-dimensional spaces of
analysis.

%%%%%%%%%%%%%%%%%%%%%%%%%%%%%%%%%%%%%%%%%%%%%%%%%%%%%%%%%%%%
\subsection{\(L^p\) Spaces Revisited}
\label{subsec:lp-functional-analysis}
%%%%%%%%%%%%%%%%%%%%%%%%%%%%%%%%%%%%%%%%%%%%%%%%%%%%%%%%%%%%

From Lebesgue integration,

\[
L^p(X,\mu)
\]

consists of measurable functions modulo equality almost everywhere for
which

\[
\int_X|f|^p\,d\mu<\infty.
\]

For

\[
1\leq p<\infty,
\]

the norm is

\[
\boxed{
\|f\|_p
=
\left(
\int_X|f|^p\,d\mu
\right)^{1/p}.
}
\]

For \(p=\infty\),

\[
\boxed{
\|f\|_\infty
=
\operatorname*{ess\,sup}|f|.
}
\]

These spaces form some of the most important examples in functional
analysis.

%%%%%%%%%%%%%%%%%%%%%%%%%%%%%%%%%%%%%%%%%%%%%%%%%%%%%%%%%%%%
\subsection{Banach Spaces}
\label{subsec:banach-spaces}
%%%%%%%%%%%%%%%%%%%%%%%%%%%%%%%%%%%%%%%%%%%%%%%%%%%%%%%%%%%%

\begin{definitionbox}{Banach Space}
A Banach space is a normed vector space that is complete with respect to
the metric induced by its norm.
\end{definitionbox}

Thus every Cauchy sequence in a Banach space converges to an element of
the same space.

Symbolically,

\[
\boxed{
(x_n)\text{ Cauchy in }X
\Longrightarrow
x_n\to x
\text{ for some }x\in X.
}
\]

Completeness is essential because limiting processes are central to
analysis.

%%%%%%%%%%%%%%%%%%%%%%%%%%%%%%%%%%%%%%%%%%%%%%%%%%%%%%%%%%%%
\subsection{Examples of Banach Spaces}
\label{subsec:examples-banach-spaces}
%%%%%%%%%%%%%%%%%%%%%%%%%%%%%%%%%%%%%%%%%%%%%%%%%%%%%%%%%%%%

Important Banach spaces include

\[
\boxed{
\R^n,
\qquad
\mathbb{C}^n,
}
\]

\[
\boxed{
\ell^p,
\qquad
1\leq p\leq\infty,
}
\]

\[
\boxed{
L^p(X,\mu),
\qquad
1\leq p\leq\infty,
}
\]

and

\[
\boxed{
C(K)
}
\]

with the supremum norm when \(K\) is compact.

%%%%%%%%%%%%%%%%%%%%%%%%%%%%%%%%%%%%%%%%%%%%%%%%%%%%%%%%%%%%
\subsection{Incomplete Normed Spaces}
\label{subsec:incomplete-normed-spaces}
%%%%%%%%%%%%%%%%%%%%%%%%%%%%%%%%%%%%%%%%%%%%%%%%%%%%%%%%%%%%

Not every normed space is complete.

For example, let

\[
P([0,1])
\]

denote the vector space of polynomial functions on \([0,1]\) with the
supremum norm.

There exist sequences of polynomials that converge uniformly to
continuous functions that are not polynomials.

Therefore

\[
\boxed{
P([0,1])
}
\]

is not complete in the supremum norm.

Its completion is closely related to

\[
C([0,1]).
\]

%%%%%%%%%%%%%%%%%%%%%%%%%%%%%%%%%%%%%%%%%%%%%%%%%%%%%%%%%%%%
\subsection{Completion}
\label{subsec:completion-normed-space}
%%%%%%%%%%%%%%%%%%%%%%%%%%%%%%%%%%%%%%%%%%%%%%%%%%%%%%%%%%%%

Every normed space \(X\) can be embedded densely into a Banach space
\(\widehat{X}\), called its completion.

The construction generalizes the passage

\[
\Q\longrightarrow\R.
\]

The completion adds limits of Cauchy sequences that are missing from the
original space.

%%%%%%%%%%%%%%%%%%%%%%%%%%%%%%%%%%%%%%%%%%%%%%%%%%%%%%%%%%%%
\subsection{Linear Operators}
\label{subsec:linear-operators-functional}
%%%%%%%%%%%%%%%%%%%%%%%%%%%%%%%%%%%%%%%%%%%%%%%%%%%%%%%%%%%%

Let \(X\) and \(Y\) be vector spaces.

A map

\[
T:X\to Y
\]

is linear if

\[
\boxed{
T(\alpha x+\beta y)
=
\alpha T(x)+\beta T(y)
}
\]

for all vectors \(x,y\) and scalars \(\alpha,\beta\).

In finite-dimensional linear algebra, linear transformations are
represented by matrices.

In functional analysis, operators may act on functions or infinite
sequences.

%%%%%%%%%%%%%%%%%%%%%%%%%%%%%%%%%%%%%%%%%%%%%%%%%%%%%%%%%%%%
\subsection{Bounded Linear Operators}
\label{subsec:bounded-linear-operators}
%%%%%%%%%%%%%%%%%%%%%%%%%%%%%%%%%%%%%%%%%%%%%%%%%%%%%%%%%%%%

\begin{definitionbox}{Bounded Linear Operator}
Let \(X\) and \(Y\) be normed spaces.

A linear operator

\[
T:X\to Y
\]

is bounded if there exists \(C\geq0\) such that

\[
\boxed{
\|Tx\|_Y
\leq
C\|x\|_X
}
\]

for every \(x\in X\).
\end{definitionbox}

The word \emph{bounded} here refers to bounded growth relative to the norm,
not necessarily to a bounded range.

%%%%%%%%%%%%%%%%%%%%%%%%%%%%%%%%%%%%%%%%%%%%%%%%%%%%%%%%%%%%
\subsection{Boundedness and Continuity}
\label{subsec:boundedness-continuity-linear}
%%%%%%%%%%%%%%%%%%%%%%%%%%%%%%%%%%%%%%%%%%%%%%%%%%%%%%%%%%%%

For linear operators, boundedness and continuity are equivalent.

\begin{theorem}
Let \(T:X\to Y\) be linear between normed spaces.

The following are equivalent:

\begin{enumerate}
    \item \(T\) is bounded;
    \item \(T\) is continuous everywhere;
    \item \(T\) is continuous at \(0\).
\end{enumerate}
\end{theorem}

This equivalence is one reason bounded operators are the natural
morphisms of normed-space theory.

%%%%%%%%%%%%%%%%%%%%%%%%%%%%%%%%%%%%%%%%%%%%%%%%%%%%%%%%%%%%
\subsection{Operator Norm}
\label{subsec:operator-norm}
%%%%%%%%%%%%%%%%%%%%%%%%%%%%%%%%%%%%%%%%%%%%%%%%%%%%%%%%%%%%

\begin{definitionbox}{Operator Norm}
For a bounded linear operator

\[
T:X\to Y,
\]

define

\[
\boxed{
\|T\|
=
\sup_{\|x\|_X\leq1}
\|Tx\|_Y.
}
\]
\end{definitionbox}

Equivalent formulas include

\[
\boxed{
\|T\|
=
\sup_{\|x\|_X=1}\|Tx\|_Y
}
\]

for \(X\neq\{0\}\), and

\[
\boxed{
\|T\|
=
\sup_{x\neq0}
\frac{\|Tx\|_Y}{\|x\|_X}.
}
\]

The defining inequality becomes

\[
\boxed{
\|Tx\|_Y
\leq
\|T\|\|x\|_X.
}
\]

%%%%%%%%%%%%%%%%%%%%%%%%%%%%%%%%%%%%%%%%%%%%%%%%%%%%%%%%%%%%
\subsection{Space of Bounded Operators}
\label{subsec:bounded-operator-space}
%%%%%%%%%%%%%%%%%%%%%%%%%%%%%%%%%%%%%%%%%%%%%%%%%%%%%%%%%%%%

The collection of bounded linear operators from \(X\) to \(Y\) is denoted

\[
\boxed{
\mathcal{B}(X,Y).
}
\]

If

\[
X=Y,
\]

we write

\[
\boxed{
\mathcal{B}(X).
}
\]

If \(Y\) is Banach, then

\[
\boxed{
\mathcal{B}(X,Y)
}
\]

is a Banach space under the operator norm.

%%%%%%%%%%%%%%%%%%%%%%%%%%%%%%%%%%%%%%%%%%%%%%%%%%%%%%%%%%%%
\subsection{Worked Example: Operator Norm}
\label{subsec:worked-operator-norm}
%%%%%%%%%%%%%%%%%%%%%%%%%%%%%%%%%%%%%%%%%%%%%%%%%%%%%%%%%%%%

\begin{workedexamplebox}{Multiplication Operator}

Define

\[
T:C([0,1])\to C([0,1])
\]

by

\[
(Tf)(x)=xf(x).
\]

Using the supremum norm,

\[
|(Tf)(x)|
=
|x||f(x)|
\leq
\|f\|_\infty.
\]

Therefore

\[
\|Tf\|_\infty
\leq
\|f\|_\infty,
\]

so

\[
\|T\|\leq1.
\]

Take

\[
f(x)=1.
\]

Then

\[
\|f\|_\infty=1
\]

and

\[
\|Tf\|_\infty
=
\sup_{0\leq x\leq1}x
=
1.
\]

Hence

\[
\begin{answerbox}
\[
\boxed{
\|T\|=1.
}
\]
\end{answerbox}
\end{workedexamplebox}

%%%%%%%%%%%%%%%%%%%%%%%%%%%%%%%%%%%%%%%%%%%%%%%%%%%%%%%%%%%%
\subsection{Unbounded Operators}
\label{subsec:unbounded-operators}
%%%%%%%%%%%%%%%%%%%%%%%%%%%%%%%%%%%%%%%%%%%%%%%%%%%%%%%%%%%%

Not every important linear operator is bounded.

Consider differentiation

\[
D:C^1([0,1])\to C([0,1]),
\qquad
Df=f'.
\]

If both spaces are equipped only with the supremum norm

\[
\|f\|_\infty,
\]

then differentiation is not bounded.

For example,

\[
f_n(x)=\frac{\sin(nx)}{n}
\]

satisfies

\[
\|f_n\|_\infty\leq\frac1n,
\]

while

\[
f_n'(x)=\cos(nx)
\]

satisfies

\[
\|f_n'\|_\infty=1.
\]

Thus no constant \(C\) can satisfy

\[
\|f'\|_\infty
\leq
C\|f\|_\infty
\]

for every \(f\in C^1([0,1])\).

Unbounded operators become especially important in differential equations
and quantum mechanics.

%%%%%%%%%%%%%%%%%%%%%%%%%%%%%%%%%%%%%%%%%%%%%%%%%%%%%%%%%%%%
\subsection{Linear Functionals}
\label{subsec:linear-functionals}
%%%%%%%%%%%%%%%%%%%%%%%%%%%%%%%%%%%%%%%%%%%%%%%%%%%%%%%%%%%%

A linear functional is a linear map from a vector space to its scalar
field.

Thus

\[
\boxed{
\varphi:X\to\R
}
\]

or

\[
\boxed{
\varphi:X\to\mathbb{C}.
}
\]

The lowercase Greek letter \(\varphi\) is \emph{phi} or \emph{varphi} and
is commonly used for functionals.

For example,

\[
\boxed{
\varphi(f)
=
\int_0^1f(x)\,dx
}
\]

defines a linear functional on \(C([0,1])\).

%%%%%%%%%%%%%%%%%%%%%%%%%%%%%%%%%%%%%%%%%%%%%%%%%%%%%%%%%%%%
\subsection{Continuous Dual Space}
\label{subsec:continuous-dual-space}
%%%%%%%%%%%%%%%%%%%%%%%%%%%%%%%%%%%%%%%%%%%%%%%%%%%%%%%%%%%%

\begin{definitionbox}{Continuous Dual Space}
The continuous dual of a normed space \(X\), denoted

\[
\boxed{
X^*,
}
\]

is the space of all bounded linear functionals on \(X\).
\end{definitionbox}

Thus

\[
\boxed{
X^*
=
\mathcal{B}(X,\mathbb{F}),
}
\]

where

\[
\mathbb{F}
=
\R
\quad\text{or}\quad
\mathbb{C}.
\]

The dual norm is

\[
\boxed{
\|\varphi\|
=
\sup_{\|x\|\leq1}
|\varphi(x)|.
}
\]

%%%%%%%%%%%%%%%%%%%%%%%%%%%%%%%%%%%%%%%%%%%%%%%%%%%%%%%%%%%%
\subsection{Worked Example: Integral Functional}
\label{subsec:worked-integral-functional}
%%%%%%%%%%%%%%%%%%%%%%%%%%%%%%%%%%%%%%%%%%%%%%%%%%%%%%%%%%%%

\begin{workedexamplebox}{A Functional on \(C([0,1])\)}

Define

\[
\varphi(f)
=
\int_0^1f(x)\,dx.
\]

Then

\[
|\varphi(f)|
\leq
\int_0^1|f(x)|\,dx
\leq
\|f\|_\infty.
\]

Therefore

\[
\|\varphi\|\leq1.
\]

For

\[
f(x)=1,
\]

we have

\[
\|f\|_\infty=1
\]

and

\[
\varphi(f)=1.
\]

Hence

\[
\begin{answerbox}
\[
\boxed{
\|\varphi\|=1.
}
\]
\end{answerbox}
\end{workedexamplebox}

%%%%%%%%%%%%%%%%%%%%%%%%%%%%%%%%%%%%%%%%%%%%%%%%%%%%%%%%%%%%
\subsection{Hahn--Banach Theorem}
\label{subsec:hahn-banach-theorem}
%%%%%%%%%%%%%%%%%%%%%%%%%%%%%%%%%%%%%%%%%%%%%%%%%%%%%%%%%%%%

One of the foundational results of functional analysis is the
Hahn--Banach Theorem.

\begin{theorem}[Hahn--Banach Theorem, Normed-Space Form]
Let \(M\) be a linear subspace of a normed vector space \(X\), and let

\[
\varphi:M\to\mathbb{F}
\]

be a bounded linear functional.

Then there exists a bounded linear functional

\[
\widetilde{\varphi}:X\to\mathbb{F}
\]

such that

\[
\boxed{
\widetilde{\varphi}|_M=\varphi
}
\]

and

\[
\boxed{
\|\widetilde{\varphi}\|
=
\|\varphi\|.
}
\]
\end{theorem}

Thus continuous linear functionals can be extended from subspaces without
increasing their norm.

%%%%%%%%%%%%%%%%%%%%%%%%%%%%%%%%%%%%%%%%%%%%%%%%%%%%%%%%%%%%
\subsection{Consequences of Hahn--Banach}
\label{subsec:hahn-banach-consequences}
%%%%%%%%%%%%%%%%%%%%%%%%%%%%%%%%%%%%%%%%%%%%%%%%%%%%%%%%%%%%

Hahn--Banach implies that continuous linear functionals are abundant.

For every nonzero \(x\in X\), there exists

\[
\varphi\in X^*
\]

such that

\[
\boxed{
\|\varphi\|=1
}
\]

and

\[
\boxed{
\varphi(x)=\|x\|.
}
\]

Consequently, dual functionals can distinguish points:

\[
x\neq y
\]

implies there exists \(\varphi\in X^*\) such that

\[
\boxed{
\varphi(x)\neq\varphi(y).
}
\]

%%%%%%%%%%%%%%%%%%%%%%%%%%%%%%%%%%%%%%%%%%%%%%%%%%%%%%%%%%%%
\subsection{Canonical Embedding into the Double Dual}
\label{subsec:double-dual-embedding}
%%%%%%%%%%%%%%%%%%%%%%%%%%%%%%%%%%%%%%%%%%%%%%%%%%%%%%%%%%%%

The dual of \(X^*\) is

\[
X^{**}.
\]

Each \(x\in X\) determines a functional on \(X^*\) by

\[
\boxed{
J(x)(\varphi)
=
\varphi(x).
}
\]

Thus there is a canonical map

\[
\boxed{
J:X\to X^{**}.
}
\]

Hahn--Banach implies that this map is an isometry:

\[
\boxed{
\|J(x)\|=\|x\|.
}
\]

%%%%%%%%%%%%%%%%%%%%%%%%%%%%%%%%%%%%%%%%%%%%%%%%%%%%%%%%%%%%
\subsection{Reflexive Spaces}
\label{subsec:reflexive-spaces}
%%%%%%%%%%%%%%%%%%%%%%%%%%%%%%%%%%%%%%%%%%%%%%%%%%%%%%%%%%%%

\begin{definitionbox}{Reflexive Space}
A Banach space \(X\) is reflexive if the canonical embedding

\[
J:X\to X^{**}
\]

is surjective.
\end{definitionbox}

Thus

\[
\boxed{
X\cong X^{**}.
}
\]

Important examples include

\[
\boxed{
L^p,
\qquad
\ell^p,
\qquad
1<p<\infty.
}
\]

Every Hilbert space is reflexive.

In contrast,

\[
L^1
\]

and

\[
\ell^1
\]

are generally not reflexive.

%%%%%%%%%%%%%%%%%%%%%%%%%%%%%%%%%%%%%%%%%%%%%%%%%%%%%%%%%%%%
\subsection{Uniform Boundedness Principle}
\label{subsec:uniform-boundedness-principle}
%%%%%%%%%%%%%%%%%%%%%%%%%%%%%%%%%%%%%%%%%%%%%%%%%%%%%%%%%%%%

The Uniform Boundedness Principle is also called the Banach--Steinhaus
Theorem.

\begin{theorem}[Uniform Boundedness Principle]
Let \(X\) be a Banach space and \(Y\) a normed space.

Suppose

\[
\mathcal{F}\subseteq\mathcal{B}(X,Y)
\]

is a family of bounded linear operators such that for every \(x\in X\),

\[
\boxed{
\sup_{T\in\mathcal{F}}\|Tx\|<\infty.
}
\]

Then

\[
\boxed{
\sup_{T\in\mathcal{F}}\|T\|<\infty.
}
\]
\end{theorem}

Thus pointwise boundedness of a family of operators on a Banach space
forces uniform boundedness in operator norm.

%%%%%%%%%%%%%%%%%%%%%%%%%%%%%%%%%%%%%%%%%%%%%%%%%%%%%%%%%%%%
\subsection{Open Mapping Theorem}
\label{subsec:open-mapping-functional}
%%%%%%%%%%%%%%%%%%%%%%%%%%%%%%%%%%%%%%%%%%%%%%%%%%%%%%%%%%%%

\begin{theorem}[Open Mapping Theorem]
Let \(X\) and \(Y\) be Banach spaces.

If

\[
T:X\to Y
\]

is bounded, linear, and surjective, then \(T\) is an open map.
\end{theorem}

That is, \(T\) maps open subsets of \(X\) to open subsets of \(Y\).

A major consequence is the bounded inverse theorem.

%%%%%%%%%%%%%%%%%%%%%%%%%%%%%%%%%%%%%%%%%%%%%%%%%%%%%%%%%%%%
\subsection{Bounded Inverse Theorem}
\label{subsec:bounded-inverse-theorem}
%%%%%%%%%%%%%%%%%%%%%%%%%%%%%%%%%%%%%%%%%%%%%%%%%%%%%%%%%%%%

\begin{theorem}[Bounded Inverse Theorem]
Let \(X\) and \(Y\) be Banach spaces.

If

\[
T:X\to Y
\]

is bounded, linear, and bijective, then

\[
\boxed{
T^{-1}:Y\to X
}
\]

is also bounded.
\end{theorem}

Thus an algebraically invertible bounded operator between Banach spaces
automatically has a continuous inverse.

%%%%%%%%%%%%%%%%%%%%%%%%%%%%%%%%%%%%%%%%%%%%%%%%%%%%%%%%%%%%
\subsection{Closed Operators}
\label{subsec:closed-operators}
%%%%%%%%%%%%%%%%%%%%%%%%%%%%%%%%%%%%%%%%%%%%%%%%%%%%%%%%%%%%

The graph of an operator \(T:X\to Y\) is

\[
\boxed{
\operatorname{Graph}(T)
=
\{(x,Tx):x\in X\}.
}
\]

The operator is called closed if its graph is closed in

\[
X\times Y.
\]

Equivalently, if

\[
x_n\to x
\]

and

\[
Tx_n\to y,
\]

then closedness requires

\[
\boxed{
Tx=y.
}
\]

%%%%%%%%%%%%%%%%%%%%%%%%%%%%%%%%%%%%%%%%%%%%%%%%%%%%%%%%%%%%
\subsection{Closed Graph Theorem}
\label{subsec:closed-graph-theorem}
%%%%%%%%%%%%%%%%%%%%%%%%%%%%%%%%%%%%%%%%%%%%%%%%%%%%%%%%%%%%

\begin{theorem}[Closed Graph Theorem]
Let \(X\) and \(Y\) be Banach spaces.

A linear operator

\[
T:X\to Y
\]

defined on all of \(X\) is bounded if and only if its graph is closed.
\end{theorem}

This provides a powerful method for proving continuity without directly
estimating the operator norm.

%%%%%%%%%%%%%%%%%%%%%%%%%%%%%%%%%%%%%%%%%%%%%%%%%%%%%%%%%%%%
\subsection{The Three Fundamental Banach-Space Theorems}
\label{subsec:three-banach-theorems}
%%%%%%%%%%%%%%%%%%%%%%%%%%%%%%%%%%%%%%%%%%%%%%%%%%%%%%%%%%%%

Three major results form a central structural core of functional analysis:

\[
\boxed{
\begin{array}{c}
\text{Uniform Boundedness Principle}\\[2mm]
\text{Open Mapping Theorem}\\[2mm]
\text{Closed Graph Theorem}
\end{array}
}
\]

All depend critically on completeness.

Their proofs ultimately rely on the Baire Category Theorem from topology.

%%%%%%%%%%%%%%%%%%%%%%%%%%%%%%%%%%%%%%%%%%%%%%%%%%%%%%%%%%%%
\subsection{Inner-Product Spaces}
\label{subsec:inner-product-spaces}
%%%%%%%%%%%%%%%%%%%%%%%%%%%%%%%%%%%%%%%%%%%%%%%%%%%%%%%%%%%%

\begin{definitionbox}{Inner Product}
An inner product on a complex vector space \(H\) is a function

\[
\langle\cdot,\cdot\rangle:H\times H\to\mathbb{C}
\]

satisfying, under the convention used here,

\[
\boxed{
\langle \alpha x+\beta y,z\rangle
=
\alpha\langle x,z\rangle
+
\beta\langle y,z\rangle,
}
\]

\[
\boxed{
\langle x,y\rangle
=
\overline{\langle y,x\rangle},
}
\]

and

\[
\boxed{
\langle x,x\rangle>0
\qquad
(x\neq0).
}
\]
\end{definitionbox}

This convention takes the inner product to be linear in the first
argument and conjugate-linear in the second.

For real vector spaces, conjugation disappears.

%%%%%%%%%%%%%%%%%%%%%%%%%%%%%%%%%%%%%%%%%%%%%%%%%%%%%%%%%%%%
\subsection{Norm Induced by an Inner Product}
\label{subsec:inner-product-norm}
%%%%%%%%%%%%%%%%%%%%%%%%%%%%%%%%%%%%%%%%%%%%%%%%%%%%%%%%%%%%

Every inner product determines a norm:

\[
\boxed{
\|x\|
=
\sqrt{\langle x,x\rangle}.
}
\]

Not every norm comes from an inner product.

Norms induced by inner products satisfy the parallelogram law

\[
\boxed{
\|x+y\|^2+\|x-y\|^2
=
2\|x\|^2+2\|y\|^2.
}
\]

%%%%%%%%%%%%%%%%%%%%%%%%%%%%%%%%%%%%%%%%%%%%%%%%%%%%%%%%%%%%
\subsection{Cauchy--Schwarz Inequality}
\label{subsec:cauchy-schwarz-functional}
%%%%%%%%%%%%%%%%%%%%%%%%%%%%%%%%%%%%%%%%%%%%%%%%%%%%%%%%%%%%

\begin{theorem}[Cauchy--Schwarz Inequality]
For every \(x,y\) in an inner-product space,

\[
\boxed{
|\langle x,y\rangle|
\leq
\|x\|\|y\|.
}
\]
\end{theorem}

Equality holds precisely when \(x\) and \(y\) are linearly dependent,
assuming neither is zero.

The Cauchy--Schwarz inequality implies the triangle inequality for the
induced norm.

%%%%%%%%%%%%%%%%%%%%%%%%%%%%%%%%%%%%%%%%%%%%%%%%%%%%%%%%%%%%
\subsection{Hilbert Spaces}
\label{subsec:hilbert-spaces}
%%%%%%%%%%%%%%%%%%%%%%%%%%%%%%%%%%%%%%%%%%%%%%%%%%%%%%%%%%%%

\begin{definitionbox}{Hilbert Space}
A Hilbert space is an inner-product space that is complete with respect to
the norm induced by its inner product.
\end{definitionbox}

Thus a Hilbert space combines

\[
\boxed{
\text{linear structure}
+
\text{geometry}
+
\text{completeness}.
}
\]

Important examples include

\[
\boxed{
\mathbb{C}^n,
\qquad
\ell^2,
\qquad
L^2(X,\mu).
}
\]

%%%%%%%%%%%%%%%%%%%%%%%%%%%%%%%%%%%%%%%%%%%%%%%%%%%%%%%%%%%%
\subsection{Orthogonality}
\label{subsec:orthogonality-hilbert}
%%%%%%%%%%%%%%%%%%%%%%%%%%%%%%%%%%%%%%%%%%%%%%%%%%%%%%%%%%%%

Vectors \(x,y\in H\) are orthogonal if

\[
\boxed{
\langle x,y\rangle=0.
}
\]

This is written

\[
\boxed{
x\perp y.
}
\]

If \(x\perp y\), then the Pythagorean identity holds:

\[
\boxed{
\|x+y\|^2
=
\|x\|^2+\|y\|^2.
}
\]

%%%%%%%%%%%%%%%%%%%%%%%%%%%%%%%%%%%%%%%%%%%%%%%%%%%%%%%%%%%%
\subsection{Orthogonal Complements}
\label{subsec:orthogonal-complements}
%%%%%%%%%%%%%%%%%%%%%%%%%%%%%%%%%%%%%%%%%%%%%%%%%%%%%%%%%%%%

For a subset \(M\subseteq H\), define

\[
\boxed{
M^\perp
=
\{
x\in H:
\langle x,m\rangle=0
\text{ for every }m\in M
\}.
}
\]

The set \(M^\perp\) is always a closed linear subspace.

For a closed subspace \(M\) of a Hilbert space,

\[
\boxed{
H=M\oplus M^\perp.
}
\]

%%%%%%%%%%%%%%%%%%%%%%%%%%%%%%%%%%%%%%%%%%%%%%%%%%%%%%%%%%%%
\subsection{Projection Theorem}
\label{subsec:projection-theorem}
%%%%%%%%%%%%%%%%%%%%%%%%%%%%%%%%%%%%%%%%%%%%%%%%%%%%%%%%%%%%

\begin{theorem}[Projection Theorem]
Let \(M\) be a closed linear subspace of a Hilbert space \(H\).

For every \(x\in H\), there exists a unique \(m\in M\) such that

\[
\boxed{
\|x-m\|
=
\inf_{y\in M}\|x-y\|.
}
\]

Moreover,

\[
\boxed{
x-m\in M^\perp.
}
\]
\end{theorem}

Thus every vector has a unique decomposition

\[
\boxed{
x=P_Mx+(x-P_Mx),
}
\]

where

\[
P_Mx\in M
\]

and

\[
x-P_Mx\in M^\perp.
\]

The operator \(P_M\) is the orthogonal projection onto \(M\).

%%%%%%%%%%%%%%%%%%%%%%%%%%%%%%%%%%%%%%%%%%%%%%%%%%%%%%%%%%%%
\subsection{Worked Example: Orthogonal Projection}
\label{subsec:worked-orthogonal-projection}
%%%%%%%%%%%%%%%%%%%%%%%%%%%%%%%%%%%%%%%%%%%%%%%%%%%%%%%%%%%%

\begin{workedexamplebox}{Projection onto a One-Dimensional Subspace}

Let \(u\neq0\) in a Hilbert space and let

\[
M=\operatorname{span}\{u\}.
\]

The projection of \(x\) onto \(M\) is

\[
P_Mx
=
\frac{\langle x,u\rangle}
{\langle u,u\rangle}
u.
\]

Since

\[
\langle u,u\rangle=\|u\|^2,
\]

we obtain

\[
\begin{answerbox}
\[
\boxed{
P_Mx
=
\frac{\langle x,u\rangle}
{\|u\|^2}u.
}
\]
\end{answerbox}
\end{workedexamplebox}

%%%%%%%%%%%%%%%%%%%%%%%%%%%%%%%%%%%%%%%%%%%%%%%%%%%%%%%%%%%%
\subsection{Orthonormal Sets}
\label{subsec:orthonormal-sets}
%%%%%%%%%%%%%%%%%%%%%%%%%%%%%%%%%%%%%%%%%%%%%%%%%%%%%%%%%%%%

A set

\[
\{e_\alpha\}_{\alpha\in A}
\]

is orthonormal if

\[
\boxed{
\langle e_\alpha,e_\beta\rangle
=
\begin{cases}
1, & \alpha=\beta,\\
0, & \alpha\neq\beta.
\end{cases}
}
\]

For a sequence \((e_n)\), this becomes

\[
\boxed{
\langle e_n,e_m\rangle
=
\delta_{nm},
}
\]

where \(\delta_{nm}\) is the Kronecker delta.

%%%%%%%%%%%%%%%%%%%%%%%%%%%%%%%%%%%%%%%%%%%%%%%%%%%%%%%%%%%%
\subsection{Bessel's Inequality}
\label{subsec:bessels-inequality}
%%%%%%%%%%%%%%%%%%%%%%%%%%%%%%%%%%%%%%%%%%%%%%%%%%%%%%%%%%%%

\begin{theorem}[Bessel's Inequality]
If \((e_n)\) is an orthonormal sequence in a Hilbert space \(H\), then

\[
\boxed{
\sum_{n=1}^{\infty}
|\langle x,e_n\rangle|^2
\leq
\|x\|^2
}
\]

for every \(x\in H\).
\end{theorem}

The coefficients

\[
\langle x,e_n\rangle
\]

measure the components of \(x\) in the directions \(e_n\).

%%%%%%%%%%%%%%%%%%%%%%%%%%%%%%%%%%%%%%%%%%%%%%%%%%%%%%%%%%%%
\subsection{Orthonormal Bases}
\label{subsec:orthonormal-bases}
%%%%%%%%%%%%%%%%%%%%%%%%%%%%%%%%%%%%%%%%%%%%%%%%%%%%%%%%%%%%

An orthonormal set is an orthonormal basis if its closed linear span is the
entire Hilbert space.

For a countable orthonormal basis \((e_n)\),

\[
\boxed{
x
=
\sum_{n=1}^{\infty}
\langle x,e_n\rangle e_n,
}
\]

where the series converges in the Hilbert-space norm.

This is the infinite-dimensional analogue of expressing a vector in
orthonormal coordinates.

%%%%%%%%%%%%%%%%%%%%%%%%%%%%%%%%%%%%%%%%%%%%%%%%%%%%%%%%%%%%
\subsection{Parseval's Identity}
\label{subsec:parsevals-identity-functional}
%%%%%%%%%%%%%%%%%%%%%%%%%%%%%%%%%%%%%%%%%%%%%%%%%%%%%%%%%%%%

If \((e_n)\) is an orthonormal basis of \(H\), then

\[
\boxed{
\|x\|^2
=
\sum_{n=1}^{\infty}
|\langle x,e_n\rangle|^2.
}
\]

This is Parseval's identity.

It strengthens Bessel's inequality from

\[
\leq
\]

to equality when the orthonormal system is complete.

Fourier series are a central example of this principle.

%%%%%%%%%%%%%%%%%%%%%%%%%%%%%%%%%%%%%%%%%%%%%%%%%%%%%%%%%%%%
\subsection{Gram--Schmidt Process}
\label{subsec:gram-schmidt-functional}
%%%%%%%%%%%%%%%%%%%%%%%%%%%%%%%%%%%%%%%%%%%%%%%%%%%%%%%%%%%%

Given linearly independent vectors

\[
x_1,x_2,\ldots,
\]

the Gram--Schmidt process constructs orthonormal vectors spanning the same
successive subspaces.

For example,

\[
e_1
=
\frac{x_1}{\|x_1\|},
\]

and

\[
u_2
=
x_2-\langle x_2,e_1\rangle e_1,
\]

followed by

\[
e_2
=
\frac{u_2}{\|u_2\|}.
\]

The process continues by subtracting projections onto previously
constructed orthonormal directions.

%%%%%%%%%%%%%%%%%%%%%%%%%%%%%%%%%%%%%%%%%%%%%%%%%%%%%%%%%%%%
\subsection{Riesz Representation Theorem}
\label{subsec:riesz-representation-hilbert}
%%%%%%%%%%%%%%%%%%%%%%%%%%%%%%%%%%%%%%%%%%%%%%%%%%%%%%%%%%%%

\begin{theorem}[Riesz Representation Theorem for Hilbert Spaces]
Let \(H\) be a Hilbert space.

For every bounded linear functional

\[
\varphi\in H^*,
\]

there exists a unique \(y\in H\) such that

\[
\boxed{
\varphi(x)
=
\langle x,y\rangle
}
\]

for every \(x\in H\).

Moreover,

\[
\boxed{
\|\varphi\|
=
\|y\|.
}
\]
\end{theorem}

Thus the continuous dual of a Hilbert space can be represented by the
Hilbert space itself, with the usual conjugate-linearity caveat in the
complex case.

%%%%%%%%%%%%%%%%%%%%%%%%%%%%%%%%%%%%%%%%%%%%%%%%%%%%%%%%%%%%
\subsection{Strong Convergence}
\label{subsec:strong-convergence}
%%%%%%%%%%%%%%%%%%%%%%%%%%%%%%%%%%%%%%%%%%%%%%%%%%%%%%%%%%%%

A sequence \((x_n)\) converges strongly to \(x\) if it converges in norm:

\[
\boxed{
\|x_n-x\|\to0.
}
\]

This is written

\[
\boxed{
x_n\to x.
}
\]

Strong convergence directly controls the distance between vectors.

%%%%%%%%%%%%%%%%%%%%%%%%%%%%%%%%%%%%%%%%%%%%%%%%%%%%%%%%%%%%
\subsection{Weak Convergence}
\label{subsec:weak-convergence}
%%%%%%%%%%%%%%%%%%%%%%%%%%%%%%%%%%%%%%%%%%%%%%%%%%%%%%%%%%%%

\begin{definitionbox}{Weak Convergence}
A sequence \((x_n)\) in a normed space \(X\) converges weakly to \(x\) if

\[
\boxed{
\varphi(x_n)\to\varphi(x)
}
\]

for every

\[
\varphi\in X^*.
\]
\end{definitionbox}

Weak convergence is commonly written

\[
\boxed{
x_n\rightharpoonup x.
}
\]

The symbol \(\rightharpoonup\) is read ``converges weakly to.''

%%%%%%%%%%%%%%%%%%%%%%%%%%%%%%%%%%%%%%%%%%%%%%%%%%%%%%%%%%%%
\subsection{Strong Versus Weak Convergence}
\label{subsec:strong-versus-weak}
%%%%%%%%%%%%%%%%%%%%%%%%%%%%%%%%%%%%%%%%%%%%%%%%%%%%%%%%%%%%

Strong convergence implies weak convergence:

\[
\boxed{
x_n\to x
\Longrightarrow
x_n\rightharpoonup x.
}
\]

The converse generally fails.

For example, let \((e_n)\) be the standard orthonormal basis of
\(\ell^2\).

Then

\[
\|e_n\|_2=1
\]

for every \(n\), so \(e_n\) does not converge strongly to zero.

However,

\[
\boxed{
e_n\rightharpoonup0.
}
\]

Thus weak convergence can exist even when norms do not converge to zero.

%%%%%%%%%%%%%%%%%%%%%%%%%%%%%%%%%%%%%%%%%%%%%%%%%%%%%%%%%%%%
\subsection{Weak-Star Convergence}
\label{subsec:weak-star-convergence}
%%%%%%%%%%%%%%%%%%%%%%%%%%%%%%%%%%%%%%%%%%%%%%%%%%%%%%%%%%%%

A sequence or net \((\varphi_\alpha)\) in \(X^*\) converges weak-star to
\(\varphi\in X^*\) if

\[
\boxed{
\varphi_\alpha(x)
\to
\varphi(x)
}
\]

for every \(x\in X\).

This topology is denoted

\[
\boxed{
\sigma(X^*,X).
}
\]

Weak-star compactness becomes important in optimization, PDEs, measure
theory, and the calculus of variations.

%%%%%%%%%%%%%%%%%%%%%%%%%%%%%%%%%%%%%%%%%%%%%%%%%%%%%%%%%%%%
\subsection{Banach--Alaoglu Theorem}
\label{subsec:banach-alaoglu}
%%%%%%%%%%%%%%%%%%%%%%%%%%%%%%%%%%%%%%%%%%%%%%%%%%%%%%%%%%%%

\begin{theorem}[Banach--Alaoglu Theorem]
The closed unit ball

\[
\boxed{
B_{X^*}
=
\{\varphi\in X^*:\|\varphi\|\leq1\}
}
\]

is compact in the weak-star topology.
\end{theorem}

In infinite-dimensional spaces, norm-closed bounded sets generally need
not be norm compact.

Banach--Alaoglu recovers an important form of compactness by weakening the
topology.

%%%%%%%%%%%%%%%%%%%%%%%%%%%%%%%%%%%%%%%%%%%%%%%%%%%%%%%%%%%%
\subsection{Compactness in Infinite Dimensions}
\label{subsec:compactness-infinite-dimensional}
%%%%%%%%%%%%%%%%%%%%%%%%%%%%%%%%%%%%%%%%%%%%%%%%%%%%%%%%%%%%

In finite-dimensional normed spaces, closed and bounded sets are compact.

This fails in infinite dimensions.

For example, the standard basis

\[
(e_n)
\]

of \(\ell^2\) lies in the closed unit ball, but

\[
\|e_n-e_m\|_2
=
\sqrt{2}
\qquad
(n\neq m).
\]

Thus it has no norm-convergent subsequence.

Therefore the closed unit ball of \(\ell^2\) is not norm compact.

%%%%%%%%%%%%%%%%%%%%%%%%%%%%%%%%%%%%%%%%%%%%%%%%%%%%%%%%%%%%
\subsection{Compact Operators}
\label{subsec:compact-operators}
%%%%%%%%%%%%%%%%%%%%%%%%%%%%%%%%%%%%%%%%%%%%%%%%%%%%%%%%%%%%

\begin{definitionbox}{Compact Operator}
A bounded linear operator

\[
T:X\to Y
\]

is compact if it maps bounded subsets of \(X\) into relatively compact
subsets of \(Y\).
\end{definitionbox}

Equivalently, if \((x_n)\) is bounded in \(X\), then \((Tx_n)\) has a
convergent subsequence in \(Y\).

Every finite-rank bounded operator is compact.

%%%%%%%%%%%%%%%%%%%%%%%%%%%%%%%%%%%%%%%%%%%%%%%%%%%%%%%%%%%%
\subsection{Finite-Rank Operators}
\label{subsec:finite-rank-operators}
%%%%%%%%%%%%%%%%%%%%%%%%%%%%%%%%%%%%%%%%%%%%%%%%%%%%%%%%%%%%

An operator \(T:X\to Y\) has finite rank if

\[
\boxed{
\dim(\operatorname{Ran}T)<\infty.
}
\]

For example,

\[
Tx
=
\sum_{k=1}^{n}
\varphi_k(x)y_k,
\]

where

\[
\varphi_k\in X^*
\]

and

\[
y_k\in Y,
\]

has finite-dimensional range.

Finite-rank operators are basic examples of compact operators.

%%%%%%%%%%%%%%%%%%%%%%%%%%%%%%%%%%%%%%%%%%%%%%%%%%%%%%%%%%%%
\subsection{Integral Operators}
\label{subsec:integral-operators-functional}
%%%%%%%%%%%%%%%%%%%%%%%%%%%%%%%%%%%%%%%%%%%%%%%%%%%%%%%%%%%%

Many operators arise from kernels.

For example,

\[
\boxed{
(Tf)(x)
=
\int_a^b
K(x,y)f(y)\,dy.
}
\]

The function

\[
K(x,y)
\]

is called the kernel of the operator.

Under suitable regularity assumptions, such integral operators may be
bounded or compact.

Integral operators occur throughout differential equations, probability,
quantum mechanics, and mathematical physics.

%%%%%%%%%%%%%%%%%%%%%%%%%%%%%%%%%%%%%%%%%%%%%%%%%%%%%%%%%%%%
\subsection{Eigenvalues and Eigenvectors}
\label{subsec:eigenvalues-functional}
%%%%%%%%%%%%%%%%%%%%%%%%%%%%%%%%%%%%%%%%%%%%%%%%%%%%%%%%%%%%

Let

\[
T:X\to X
\]

be linear.

A nonzero vector \(x\) is an eigenvector with eigenvalue \(\lambda\) if

\[
\boxed{
Tx=\lambda x.
}
\]

Equivalently,

\[
\boxed{
(T-\lambda I)x=0.
}
\]

Thus \(\lambda\) is an eigenvalue precisely when

\[
T-\lambda I
\]

fails to be injective.

%%%%%%%%%%%%%%%%%%%%%%%%%%%%%%%%%%%%%%%%%%%%%%%%%%%%%%%%%%%%
\subsection{Spectrum}
\label{subsec:spectrum-operator}
%%%%%%%%%%%%%%%%%%%%%%%%%%%%%%%%%%%%%%%%%%%%%%%%%%%%%%%%%%%%

In infinite-dimensional spaces, eigenvalues alone do not capture all
possible failures of invertibility.

\begin{definitionbox}{Spectrum}
Let \(X\) be a complex Banach space and

\[
T\in\mathcal{B}(X).
\]

The spectrum of \(T\) is

\[
\boxed{
\sigma(T)
=
\{
\lambda\in\mathbb{C}:
T-\lambda I
\text{ is not invertible in }\mathcal{B}(X)
\}.
}
\]
\end{definitionbox}

The lowercase Greek letter \(\sigma\) is \emph{sigma}.

%%%%%%%%%%%%%%%%%%%%%%%%%%%%%%%%%%%%%%%%%%%%%%%%%%%%%%%%%%%%
\subsection{Resolvent}
\label{subsec:resolvent-operator}
%%%%%%%%%%%%%%%%%%%%%%%%%%%%%%%%%%%%%%%%%%%%%%%%%%%%%%%%%%%%

The resolvent set is

\[
\boxed{
\rho(T)
=
\mathbb{C}\setminus\sigma(T).
}
\]

For

\[
\lambda\in\rho(T),
\]

the resolvent operator is

\[
\boxed{
R(\lambda,T)
=
(T-\lambda I)^{-1}.
}
\]

The resolvent is operator-valued and holomorphic as a function of
\(\lambda\) on \(\rho(T)\).

%%%%%%%%%%%%%%%%%%%%%%%%%%%%%%%%%%%%%%%%%%%%%%%%%%%%%%%%%%%%
\subsection{Basic Properties of the Spectrum}
\label{subsec:basic-spectrum-properties}
%%%%%%%%%%%%%%%%%%%%%%%%%%%%%%%%%%%%%%%%%%%%%%%%%%%%%%%%%%%%

If

\[
T\in\mathcal{B}(X)
\]

on a complex Banach space, then

\[
\boxed{
\sigma(T)\neq\varnothing.
}
\]

Moreover,

\[
\boxed{
\sigma(T)
\text{ is compact}
}
\]

and

\[
\boxed{
\sigma(T)
\subseteq
\{
\lambda\in\mathbb{C}:
|\lambda|\leq\|T\|
\}.
}
\]

Thus

\[
\boxed{
|\lambda|\leq\|T\|
}
\]

for every spectral value \(\lambda\).

%%%%%%%%%%%%%%%%%%%%%%%%%%%%%%%%%%%%%%%%%%%%%%%%%%%%%%%%%%%%
\subsection{Spectral Radius}
\label{subsec:spectral-radius}
%%%%%%%%%%%%%%%%%%%%%%%%%%%%%%%%%%%%%%%%%%%%%%%%%%%%%%%%%%%%

The spectral radius is

\[
\boxed{
r(T)
=
\sup_{\lambda\in\sigma(T)}
|\lambda|.
}
\]

The spectral radius formula states

\[
\boxed{
r(T)
=
\lim_{n\to\infty}
\|T^n\|^{1/n}.
}
\]

Equivalently,

\[
\boxed{
r(T)
=
\inf_{n\geq1}
\|T^n\|^{1/n}.
}
\]

This connects spectral information to asymptotic operator growth.

%%%%%%%%%%%%%%%%%%%%%%%%%%%%%%%%%%%%%%%%%%%%%%%%%%%%%%%%%%%%
\subsection{Adjoint Operators}
\label{subsec:adjoint-operators}
%%%%%%%%%%%%%%%%%%%%%%%%%%%%%%%%%%%%%%%%%%%%%%%%%%%%%%%%%%%%

Let \(H\) be a Hilbert space and

\[
T\in\mathcal{B}(H).
\]

The adjoint of \(T\) is the unique operator

\[
\boxed{
T^*:H\to H
}
\]

satisfying

\[
\boxed{
\langle Tx,y\rangle
=
\langle x,T^*y\rangle
}
\]

for all \(x,y\in H\).

The adjoint generalizes conjugate transpose from matrix theory.

%%%%%%%%%%%%%%%%%%%%%%%%%%%%%%%%%%%%%%%%%%%%%%%%%%%%%%%%%%%%
\subsection{Self-Adjoint Operators}
\label{subsec:self-adjoint-operators}
%%%%%%%%%%%%%%%%%%%%%%%%%%%%%%%%%%%%%%%%%%%%%%%%%%%%%%%%%%%%

\begin{definitionbox}{Self-Adjoint Operator}
An operator \(T\in\mathcal{B}(H)\) is self-adjoint if

\[
\boxed{
T=T^*.
}
\]
\end{definitionbox}

Self-adjoint operators generalize real symmetric and complex Hermitian
matrices.

Their spectral behavior is especially well structured.

In particular,

\[
\boxed{
\sigma(T)\subseteq\R
}
\]

for bounded self-adjoint \(T\).

%%%%%%%%%%%%%%%%%%%%%%%%%%%%%%%%%%%%%%%%%%%%%%%%%%%%%%%%%%%%
\subsection{Normal and Unitary Operators}
\label{subsec:normal-unitary-operators}
%%%%%%%%%%%%%%%%%%%%%%%%%%%%%%%%%%%%%%%%%%%%%%%%%%%%%%%%%%%%

An operator is normal if

\[
\boxed{
TT^*=T^*T.
}
\]

An operator is unitary if

\[
\boxed{
T^*T=TT^*=I.
}
\]

Unitary operators preserve inner products:

\[
\boxed{
\langle Tx,Ty\rangle
=
\langle x,y\rangle.
}
\]

Consequently,

\[
\boxed{
\|Tx\|=\|x\|.
}
\]

%%%%%%%%%%%%%%%%%%%%%%%%%%%%%%%%%%%%%%%%%%%%%%%%%%%%%%%%%%%%
\subsection{Compact Self-Adjoint Operators}
\label{subsec:compact-self-adjoint}
%%%%%%%%%%%%%%%%%%%%%%%%%%%%%%%%%%%%%%%%%%%%%%%%%%%%%%%%%%%%

Compact self-adjoint operators behave in many ways like symmetric
matrices.

\begin{theorem}[Spectral Theorem for Compact Self-Adjoint Operators]
Let \(T\) be compact and self-adjoint on a Hilbert space \(H\).

Then every nonzero spectral value is a real eigenvalue, nonzero
eigenvalues have finite multiplicity, and the only possible accumulation
point of the nonzero eigenvalues is \(0\).

Moreover, \(H\) admits an orthonormal basis adapted to the eigenspaces of
\(T\), together with the kernel when necessary.
\end{theorem}

This theorem provides an infinite-dimensional analogue of orthogonal
diagonalization.

%%%%%%%%%%%%%%%%%%%%%%%%%%%%%%%%%%%%%%%%%%%%%%%%%%%%%%%%%%%%
\subsection{Spectral Theorem Preview}
\label{subsec:spectral-theorem-preview}
%%%%%%%%%%%%%%%%%%%%%%%%%%%%%%%%%%%%%%%%%%%%%%%%%%%%%%%%%%%%

For general bounded self-adjoint operators, diagonalization may no longer
mean a discrete eigenvector expansion.

Instead, the spectral theorem represents \(T\) through a projection-valued
measure.

Symbolically,

\[
\boxed{
T
=
\int_{\sigma(T)}
\lambda\,dE(\lambda).
}
\]

This viewpoint is developed more fully in Operator Theory.

%%%%%%%%%%%%%%%%%%%%%%%%%%%%%%%%%%%%%%%%%%%%%%%%%%%%%%%%%%%%
\subsection{Functionals on \(L^p\)}
\label{subsec:duality-lp}
%%%%%%%%%%%%%%%%%%%%%%%%%%%%%%%%%%%%%%%%%%%%%%%%%%%%%%%%%%%%

For

\[
1<p<\infty
\]

and conjugate exponent \(q\) satisfying

\[
\frac1p+\frac1q=1,
\]

a function

\[
g\in L^q
\]

defines a bounded functional on \(L^p\) by

\[
\boxed{
\varphi_g(f)
=
\int_Xf\overline{g}\,d\mu.
}
\]

H\"older's inequality gives

\[
|\varphi_g(f)|
\leq
\|f\|_p\|g\|_q.
\]

Under standard hypotheses,

\[
\boxed{
(L^p)^*
\cong
L^q.
}
\]

This is one of the central duality relationships in analysis.

%%%%%%%%%%%%%%%%%%%%%%%%%%%%%%%%%%%%%%%%%%%%%%%%%%%%%%%%%%%%
\subsection{Duality of Sequence Spaces}
\label{subsec:duality-sequence-spaces}
%%%%%%%%%%%%%%%%%%%%%%%%%%%%%%%%%%%%%%%%%%%%%%%%%%%%%%%%%%%%

For

\[
1<p<\infty
\]

with conjugate exponent \(q\),

\[
\boxed{
(\ell^p)^*
\cong
\ell^q.
}
\]

In particular,

\[
\boxed{
(\ell^2)^*
\cong
\ell^2.
}
\]

Also,

\[
\boxed{
(c_0)^*
\cong
\ell^1.
}
\]

Duality is fundamental because it converts geometric and analytic
questions about vectors into questions about linear functionals.

%%%%%%%%%%%%%%%%%%%%%%%%%%%%%%%%%%%%%%%%%%%%%%%%%%%%%%%%%%%%
\subsection{Separable Spaces}
\label{subsec:separable-spaces-functional}
%%%%%%%%%%%%%%%%%%%%%%%%%%%%%%%%%%%%%%%%%%%%%%%%%%%%%%%%%%%%

\begin{definitionbox}{Separable Space}
A metric space is separable if it contains a countable dense subset.
\end{definitionbox}

Examples include

\[
\boxed{
\ell^p
\quad
(1\leq p<\infty)
}
\]

and, under standard assumptions,

\[
\boxed{
L^p
\quad
(1\leq p<\infty).
}
\]

The space

\[
\ell^\infty
\]

is not separable.

Separability often allows infinite-dimensional problems to be approximated
using countable systems.

%%%%%%%%%%%%%%%%%%%%%%%%%%%%%%%%%%%%%%%%%%%%%%%%%%%%%%%%%%%%
\subsection{The Baire Category Viewpoint}
\label{subsec:baire-functional-analysis}
%%%%%%%%%%%%%%%%%%%%%%%%%%%%%%%%%%%%%%%%%%%%%%%%%%%%%%%%%%%%

The Baire Category Theorem states, in one standard form, that a complete
metric space cannot be written as a countable union of nowhere dense
closed sets.

This topological result underlies several major theorems of functional
analysis.

Schematically,

\[
\boxed{
\text{Baire Category}
\Longrightarrow
\begin{cases}
\text{Uniform Boundedness},\\
\text{Open Mapping},\\
\text{Closed Graph}.
\end{cases}
}
\]

This is an important example of topology producing strong analytic
consequences.

%%%%%%%%%%%%%%%%%%%%%%%%%%%%%%%%%%%%%%%%%%%%%%%%%%%%%%%%%%%%
\subsection{Functional Analysis and Differential Equations}
\label{subsec:functional-analysis-differential-equations}
%%%%%%%%%%%%%%%%%%%%%%%%%%%%%%%%%%%%%%%%%%%%%%%%%%%%%%%%%%%%

A differential equation may be viewed abstractly as an operator equation

\[
\boxed{
Tu=f.
}
\]

Instead of seeking a formula for \(u\), one studies properties of the
operator \(T\):

\[
\boxed{
\text{existence},
\qquad
\text{uniqueness},
\qquad
\text{stability},
\qquad
\text{regularity}.
}
\]

This viewpoint becomes essential for partial differential equations.

%%%%%%%%%%%%%%%%%%%%%%%%%%%%%%%%%%%%%%%%%%%%%%%%%%%%%%%%%%%%
\subsection{Weak Formulations}
\label{subsec:weak-formulations-functional}
%%%%%%%%%%%%%%%%%%%%%%%%%%%%%%%%%%%%%%%%%%%%%%%%%%%%%%%%%%%%

Classical solutions may require derivatives that do not exist pointwise.

Functional analysis permits weaker notions of solution by integrating
against test functions.

Schematically, a differential equation

\[
Lu=f
\]

may be replaced by an identity such as

\[
\boxed{
a(u,v)=F(v)
}
\]

for every test function \(v\) in a suitable function space.

This leads naturally to Sobolev spaces and weak solutions.

%%%%%%%%%%%%%%%%%%%%%%%%%%%%%%%%%%%%%%%%%%%%%%%%%%%%%%%%%%%%
\subsection{Lax--Milgram Theorem Preview}
\label{subsec:lax-milgram-preview}
%%%%%%%%%%%%%%%%%%%%%%%%%%%%%%%%%%%%%%%%%%%%%%%%%%%%%%%%%%%%

A fundamental existence theorem for weak formulations is the
Lax--Milgram Theorem.

\begin{theorem}[Lax--Milgram Theorem]
Let \(H\) be a Hilbert space and let

\[
a:H\times H\to\mathbb{F}
\]

be a bounded coercive sesquilinear form.

Then for every bounded linear functional \(F\in H^*\), there exists a
unique \(u\in H\) satisfying

\[
\boxed{
a(u,v)=F(v)
}
\]

for every \(v\in H\).
\end{theorem}

This theorem is one of the major bridges from Hilbert-space theory to
partial differential equations.

%%%%%%%%%%%%%%%%%%%%%%%%%%%%%%%%%%%%%%%%%%%%%%%%%%%%%%%%%%%%
\subsection{Common Errors}
\label{subsec:functional-analysis-common-errors}
%%%%%%%%%%%%%%%%%%%%%%%%%%%%%%%%%%%%%%%%%%%%%%%%%%%%%%%%%%%%

\subsubsection{Assuming Every Normed Space Is Complete}

A normed space need not be Banach.

Completeness must be established separately.

%%%%%%%%%%%%%%%%%%%%%%%%%%%%%%%%%%%%%%%%%%%%%%%%%%%%%%%%%%%%
\subsubsection{Assuming Closed and Bounded Means Compact}

This is true in finite-dimensional normed spaces but generally false in
infinite dimensions.

%%%%%%%%%%%%%%%%%%%%%%%%%%%%%%%%%%%%%%%%%%%%%%%%%%%%%%%%%%%%
\subsubsection{Confusing Bounded Operators with Bounded Functions}

A bounded linear operator satisfies

\[
\|Tx\|
\leq
C\|x\|.
\]

Its range need not be a bounded subset.

%%%%%%%%%%%%%%%%%%%%%%%%%%%%%%%%%%%%%%%%%%%%%%%%%%%%%%%%%%%%
\subsubsection{Forgetting the Domain Norm}

Whether an operator is bounded depends on the norms placed on its domain
and codomain.

The same algebraic operator can be bounded under one choice of norms and
unbounded under another.

%%%%%%%%%%%%%%%%%%%%%%%%%%%%%%%%%%%%%%%%%%%%%%%%%%%%%%%%%%%%
\subsubsection{Treating Infinite-Dimensional Spaces Exactly Like \(\R^n\)}

Several familiar finite-dimensional results fail in infinite dimensions,
including compactness of closed bounded sets.

%%%%%%%%%%%%%%%%%%%%%%%%%%%%%%%%%%%%%%%%%%%%%%%%%%%%%%%%%%%%
\subsubsection{Confusing Strong and Weak Convergence}

Strong convergence means

\[
\|x_n-x\|\to0.
\]

Weak convergence means

\[
\varphi(x_n)\to\varphi(x)
\]

for every continuous linear functional.

%%%%%%%%%%%%%%%%%%%%%%%%%%%%%%%%%%%%%%%%%%%%%%%%%%%%%%%%%%%%
\subsubsection{Assuming Weak Convergence Implies Norm Convergence}

In general,

\[
x_n\rightharpoonup x
\]

does not imply

\[
x_n\to x
\]

in norm.

%%%%%%%%%%%%%%%%%%%%%%%%%%%%%%%%%%%%%%%%%%%%%%%%%%%%%%%%%%%%
\subsubsection{Ignoring Almost-Everywhere Equivalence in \(L^p\)}

Elements of \(L^p\) are equivalence classes of functions modulo equality
almost everywhere.

%%%%%%%%%%%%%%%%%%%%%%%%%%%%%%%%%%%%%%%%%%%%%%%%%%%%%%%%%%%%
\subsubsection{Confusing Spectrum with Eigenvalues}

In finite dimensions, the spectrum consists entirely of eigenvalues.

In infinite dimensions, spectral values may occur even when no nonzero
eigenvector exists.

%%%%%%%%%%%%%%%%%%%%%%%%%%%%%%%%%%%%%%%%%%%%%%%%%%%%%%%%%%%%
\subsubsection{Assuming Every Operator Has an Eigenbasis}

Only special classes of operators admit useful orthogonal eigenvector
decompositions.

The spectral theorem requires appropriate hypotheses.

%%%%%%%%%%%%%%%%%%%%%%%%%%%%%%%%%%%%%%%%%%%%%%%%%%%%%%%%%%%%
\subsubsection{Ignoring Completeness in Major Theorems}

The Uniform Boundedness, Open Mapping, and Closed Graph Theorems require
Banach-space hypotheses in their standard forms.

%%%%%%%%%%%%%%%%%%%%%%%%%%%%%%%%%%%%%%%%%%%%%%%%%%%%%%%%%%%%
\subsection{Applications}
\label{subsec:functional-analysis-applications}
%%%%%%%%%%%%%%%%%%%%%%%%%%%%%%%%%%%%%%%%%%%%%%%%%%%%%%%%%%%%

\begin{applicationbox}
\textbf{Partial Differential Equations.}
Banach and Hilbert spaces provide the setting for weak solutions,
existence theorems, and regularity theory.

\medskip

\textbf{Fourier Analysis.}
Fourier series and transforms are naturally interpreted using orthogonal
expansions in \(L^2\).

\medskip

\textbf{Quantum Mechanics.}
Quantum states are modeled by vectors in Hilbert spaces, while observables
are represented by operators.

\medskip

\textbf{Optimization.}
Infinite-dimensional optimization uses dual spaces, weak convergence,
convexity, and compactness.

\medskip

\textbf{Calculus of Variations.}
Minimizing sequences are frequently analyzed using weak compactness in
Banach and Hilbert spaces.

\medskip

\textbf{Integral Equations.}
Many integral equations can be written as operator equations involving
compact operators.

\medskip

\textbf{Numerical Analysis.}
Finite-dimensional approximation methods can be interpreted as
approximations of infinite-dimensional operator problems.

\medskip

\textbf{Probability Theory.}
Spaces of random variables such as \(L^1\) and \(L^2\) are Banach and
Hilbert spaces.

\medskip

\textbf{Signal Processing.}
Signals can be represented as vectors in function spaces and decomposed
using orthogonal systems.

\medskip

\textbf{Control Theory.}
Operator-theoretic methods describe infinite-dimensional dynamical
systems and stability.
\end{applicationbox}

%%%%%%%%%%%%%%%%%%%%%%%%%%%%%%%%%%%%%%%%%%%%%%%%%%%%%%%%%%%%
\subsection{Exercises}
\label{subsec:functional-analysis-exercises}
%%%%%%%%%%%%%%%%%%%%%%%%%%%%%%%%%%%%%%%%%%%%%%%%%%%%%%%%%%%%

\begin{exercise}[1]
Verify that

\[
\|x\|_1
=
\sum_{k=1}^{n}|x_k|
\]

defines a norm on \(\R^n\).
\end{exercise}

\begin{exercise}[1]
Verify that

\[
\|x\|_\infty
=
\max_k|x_k|
\]

defines a norm on \(\R^n\).
\end{exercise}

\begin{exercise}[1]
For

\[
x=(1,-2,3),
\]

compute

\[
\|x\|_1,
\qquad
\|x\|_2,
\qquad
\|x\|_\infty.
\]
\end{exercise}

\begin{exercise}[1]
Determine whether

\[
\left(\frac1n\right)_{n=1}^{\infty}
\]

belongs to

\[
\ell^1,
\qquad
\ell^2,
\qquad
\ell^\infty.
\]
\end{exercise}

\begin{exercise}[1]
Show that

\[
c_0\subseteq c\subseteq\ell^\infty.
\]
\end{exercise}

\begin{exercise}[2]
Prove that equivalent norms produce the same convergent sequences.
\end{exercise}

\begin{exercise}[2]
Show that \(c_0\) is closed in \(\ell^\infty\).
\end{exercise}

\begin{exercise}[2]
Prove that \(C([0,1])\) is complete under the supremum norm.
\end{exercise}

\begin{exercise}[2]
Define

\[
T:\ell^2\to\ell^2
\]

by

\[
T(x_1,x_2,\ldots)
=
\left(
x_1,\frac{x_2}{2},\frac{x_3}{3},\ldots
\right).
\]

Show that \(T\) is bounded and compute \(\|T\|\).
\end{exercise}

\begin{exercise}[2]
Show that the evaluation functional

\[
\varphi(f)=f(0)
\]

is bounded on \(C([0,1])\) with the supremum norm.
\end{exercise}

\begin{exercise}[2]
Compute the norm of the evaluation functional

\[
\varphi_a(f)=f(a)
\]

for fixed \(a\in[0,1]\).
\end{exercise}

\begin{exercise}[2]
Show that every finite-rank bounded operator is compact.
\end{exercise}

\begin{exercise}[2]
Let \((e_n)\) be an orthonormal sequence. Prove that

\[
\|e_n-e_m\|
=
\sqrt{2}
\]

whenever \(n\neq m\).
\end{exercise}

\begin{exercise}[2]
Show that the closed unit ball of an infinite-dimensional Hilbert space
need not be norm compact.
\end{exercise}

\begin{exercise}[3]
Prove that a linear operator is continuous if and only if it is bounded.
\end{exercise}

\begin{exercise}[3]
Prove that

\[
\|ST\|
\leq
\|S\|\|T\|
\]

for bounded operators \(T:X\to Y\) and \(S:Y\to Z\).
\end{exercise}

\begin{exercise}[3]
Use Hahn--Banach to show that for every nonzero \(x\in X\), there exists
\(\varphi\in X^*\) such that

\[
\|\varphi\|=1
\]

and

\[
\varphi(x)=\|x\|.
\]
\end{exercise}

\begin{exercise}[3]
Prove that the canonical embedding

\[
J:X\to X^{**}
\]

is an isometry.
\end{exercise}

\begin{exercise}[3]
Show that strong convergence implies weak convergence.
\end{exercise}

\begin{exercise}[3]
Prove that the standard basis \((e_n)\) of \(\ell^2\) converges weakly to
zero.
\end{exercise}

\begin{exercise}[3]
Prove Bessel's inequality.
\end{exercise}

\begin{exercise}[3]
Derive Parseval's identity from completeness of an orthonormal basis.
\end{exercise}

\begin{exercise}[3]
Show that an orthogonal projection \(P\) satisfies

\[
P^2=P.
\]
\end{exercise}

\begin{exercise}[3]
Show that if \(P\) is a nonzero orthogonal projection, then

\[
\|P\|=1.
\]
\end{exercise}

\begin{exercise}[3]
Prove that every eigenvalue of a bounded operator \(T\) satisfies

\[
|\lambda|\leq\|T\|.
\]
\end{exercise}

\begin{exercise}[3]
Show that the spectrum of a bounded self-adjoint operator is real.
\end{exercise}

\begin{exercise}[3]
Let

\[
T:\ell^2\to\ell^2
\]

be defined by

\[
T(x_1,x_2,\ldots)
=
\left(
x_1,\frac{x_2}{2},\frac{x_3}{3},\ldots
\right).
\]

Determine its eigenvalues and investigate its spectrum.
\end{exercise}

\begin{exercise}[4]
Study the proof of the Uniform Boundedness Principle using the Baire
Category Theorem.
\end{exercise}

\begin{exercise}[4]
Derive the Bounded Inverse Theorem from the Open Mapping Theorem.
\end{exercise}

\begin{exercise}[4]
Study the proof of the Closed Graph Theorem and identify exactly where
completeness is used.
\end{exercise}

\begin{exercise}[4]
Prove the Projection Theorem for closed subspaces of Hilbert spaces.
\end{exercise}

\begin{exercise}[4]
Prove the Riesz Representation Theorem for Hilbert spaces.
\end{exercise}

\begin{exercise}[4]
Show that

\[
(\ell^p)^*\cong\ell^q
\]

for

\[
1<p<\infty,
\qquad
\frac1p+\frac1q=1.
\]
\end{exercise}

\begin{exercise}[4]
Investigate why

\[
(\ell^\infty)^*
\]

is substantially larger than \(\ell^1\).
\end{exercise}

\begin{exercise}[4]
Study the proof of the Banach--Alaoglu Theorem.
\end{exercise}

\begin{exercise}[4]
Prove the spectral radius formula for a bounded operator.
\end{exercise}

\begin{exercise}[4]
Study the spectral theorem for compact self-adjoint operators and compare
it with diagonalization of real symmetric matrices.
\end{exercise}

%%%%%%%%%%%%%%%%%%%%%%%%%%%%%%%%%%%%%%%%%%%%%%%%%%%%%%%%%%%%
\subsection{Conceptual Summary}
\label{subsec:functional-analysis-summary}
%%%%%%%%%%%%%%%%%%%%%%%%%%%%%%%%%%%%%%%%%%%%%%%%%%%%%%%%%%%%

Functional analysis extends linear algebra into infinite-dimensional
spaces.

A norm provides a notion of size:

\[
\boxed{
\|x\|.
}
\]

A complete normed space is a Banach space.

Important examples are

\[
\boxed{
\ell^p,
\qquad
L^p,
\qquad
C(K).
}
\]

A linear operator is bounded when

\[
\boxed{
\|Tx\|
\leq
C\|x\|.
}
\]

For linear maps,

\[
\boxed{
\text{bounded}
\iff
\text{continuous}.
}
\]

The operator norm is

\[
\boxed{
\|T\|
=
\sup_{\|x\|\leq1}\|Tx\|.
}
\]

The continuous dual is

\[
\boxed{
X^*
=
\{\text{bounded linear functionals on }X\}.
}
\]

Hahn--Banach extends bounded linear functionals while preserving their
norm.

The Uniform Boundedness, Open Mapping, and Closed Graph Theorems reveal
the power of completeness.

An inner product provides geometry through

\[
\boxed{
\langle x,y\rangle.
}
\]

A complete inner-product space is a Hilbert space.

Orthogonal projection gives

\[
\boxed{
H=M\oplus M^\perp
}
\]

for every closed subspace \(M\).

For an orthonormal basis \((e_n)\),

\[
\boxed{
x
=
\sum_{n=1}^{\infty}
\langle x,e_n\rangle e_n
}
\]

and Parseval's identity gives

\[
\boxed{
\|x\|^2
=
\sum_{n=1}^{\infty}
|\langle x,e_n\rangle|^2.
}
\]

Weak convergence is characterized by

\[
\boxed{
x_n\rightharpoonup x
\iff
\varphi(x_n)\to\varphi(x)
\text{ for every }\varphi\in X^*.
}
\]

Compact operators generalize important finite-dimensional behavior.

The spectrum of \(T\) is

\[
\boxed{
\sigma(T)
=
\{
\lambda:
T-\lambda I
\text{ is not boundedly invertible}
\}.
}
\]

Functional analysis therefore provides the abstract language connecting
function spaces, operators, convergence, geometry, differential
equations, optimization, Fourier analysis, and mathematical physics.

%%%%%%%%%%%%%%%%%%%%%%%%%%%%%%%%%%%%%%%%%%%%%%%%%%%%%%%%%%%%
\subsection{Section Connections}
\label{subsec:functional-analysis-connections}
%%%%%%%%%%%%%%%%%%%%%%%%%%%%%%%%%%%%%%%%%%%%%%%%%%%%%%%%%%%%

\chapterconnection{
Functional Analysis builds on Linear Algebra, Topology, Real Analysis,
Measure Theory, and Lebesgue Integration by treating functions and
sequences as vectors in infinite-dimensional spaces. Banach spaces
formalize completeness, while Hilbert spaces add inner-product geometry.
The Hahn--Banach, Uniform Boundedness, Open Mapping, and Closed Graph
Theorems provide foundational tools for studying continuous linear
operators. Weak convergence and compactness become essential in the
Calculus of Variations, Optimization, and Partial Differential Equations.
Orthogonal expansions in Hilbert spaces provide the abstract foundation
for Fourier Analysis, which follows later in this chapter. Spectra,
adjoints, self-adjoint operators, and compact operators lead directly
into Operator Theory, while integral operators and weak formulations
connect the subject to Differential Equations and Mathematical Physics.
}

%%%%%%%%%%%%%%%%%%%%%%%%%%%%%%%%%%%%%%%%%%%%%%%%%%%%%%%%%%%%
% SECTION SOURCE TRACKING
%%%%%%%%%%%%%%%%%%%%%%%%%%%%%%%%%%%%%%%%%%%%%%%%%%%%%%%%%%%%

%------------------------------------------------------------
% 5.9 Functional Analysis
%------------------------------------------------------------
%
% Source basis:
%   - Standard undergraduate and beginning graduate functional analysis.
%   - Standard Banach-space and Hilbert-space theory.
%
% Used for:
%   - normed vector spaces;
%   - equivalent norms;
%   - sequence spaces;
%   - continuous-function spaces;
%   - L^p spaces;
%   - Banach spaces and completion;
%   - bounded linear operators;
%   - operator norms;
%   - continuous linear functionals;
%   - dual and double-dual spaces;
%   - Hahn--Banach Theorem;
%   - reflexive spaces;
%   - Uniform Boundedness Principle;
%   - Open Mapping Theorem;
%   - Bounded Inverse Theorem;
%   - Closed Graph Theorem;
%   - inner-product spaces;
%   - Hilbert spaces;
%   - orthogonality and projections;
%   - orthonormal bases;
%   - Bessel's inequality;
%   - Parseval's identity;
%   - Riesz Representation Theorem;
%   - strong, weak, and weak-star convergence;
%   - Banach--Alaoglu Theorem;
%   - compact operators;
%   - finite-rank and integral operators;
%   - spectrum and resolvent;
%   - spectral radius;
%   - adjoint, self-adjoint, normal, and unitary operators;
%   - compact self-adjoint spectral theory;
%   - L^p and sequence-space duality;
%   - separability;
%   - Baire-category methods;
%   - weak formulations and Lax--Milgram preview.
%
% Notes:
%   - Norms, metric spaces, completeness, and compactness are
%     developed earlier in the text.
%   - L^p spaces originate in Lebesgue Integration.
%   - Fourier Analysis develops orthogonal expansions in L^2.
%   - Operator Theory later develops spectral theory in greater depth.
%   - Weak formulations and Hilbert-space methods connect this
%     material to PDEs and the Calculus of Variations.
%
%%%%%%%%%%%%%%%%%%%%%%%%%%%%%%%%%%%%%%%%%%%%%%%%%%%%%%%%%%%%